\chapter{Python Documentation}

\section{Purpose of the Documentation}

The project includes a separate Sphinx site for the Python codebase.
Its purpose is to make the implementation inspectable at module and API level, without overloading the main report with low-level interface details.
The report and the generated site therefore have complementary roles:

\begin{itemize}
    \item the report explains the mathematical model, experimental design, software architecture, and empirical results;
    \item the Sphinx site exposes the package structure, module-level notes, function signatures, class interfaces, and docstrings.
\end{itemize}
This separation keeps the report focused on the scientific content while still making the codebase auditable.

\section{Sphinx Documentation Structure}

The Sphinx sources are stored in \path{docs/source/}.
The API reference is generated from the docstrings contained in the main Python package, \path{src/nn_option_pricing/}.

The online site contains dedicated pages for the experiment pipeline,
the additional experiments, the testing strategy, and the generated API reference.
The most important entry points are:

\begin{itemize}
    \item \href{https://github.com/matteogiorgi/nn-option-pricing/blob/main/docs/source/index.rst}{\path{index.rst}}: main documentation index.
    \item \href{https://github.com/matteogiorgi/nn-option-pricing/blob/main/docs/source/overview.rst}{\path{overview.rst}}: project orientation and repository overview.
    \item \href{https://github.com/matteogiorgi/nn-option-pricing/blob/main/docs/source/methodology.rst}{\path{methodology.rst}}: methodological description of the experiment.
    \item \href{https://github.com/matteogiorgi/nn-option-pricing/blob/main/docs/source/experiment_pipeline.rst}{\path{experiment_pipeline.rst}}: end-to-end experiment pipeline.
    \item \href{https://github.com/matteogiorgi/nn-option-pricing/blob/main/docs/source/experiments.rst}{\path{experiments.rst}}: additional experiment descriptions.
    \item \href{https://github.com/matteogiorgi/nn-option-pricing/blob/main/docs/source/testing.rst}{\path{testing.rst}}: testing and verification strategy.
    \item \href{https://github.com/matteogiorgi/nn-option-pricing/tree/main/docs/source/api}{\path{api/}}: generated API reference for the Python package.
\end{itemize}

\section{Local Documentation Build}

The HTML site can be built locally from the repository root after installing the project and Sphinx dependencies:

\begin{projectlistingbox}[shellstyle]
pip install -r requirements.txt
pip install -r requirements-docs.txt
pip install -e . --no-build-isolation
\end{projectlistingbox}
\projectlistingcaption{Installation commands for local Sphinx builds.}

The site can be built either in standard mode or in strict mode, and the generated HTML entry point can be found locally at \path{docs/build/html/index.html}:

\begin{projectlistingbox}[shellstyle]
sphinx-build -b html docs/source docs/build/html
sphinx-build -W -b html docs/source docs/build/html
\end{projectlistingbox}
\projectlistingcaption{Standard and strict Sphinx HTML build commands.}

\section{Automatic Deployment: GitHub Actions \& GitHub Pages}

The generated HTML site is not tracked in the repository.
This is intentional: \path{docs/build/html/} is a generated artifact, and keeping it out of version control avoids unnecessary file churn.
Instead, the site is built automatically through GitHub Actions and published with GitHub Pages.
The workflow is stored in \href{https://github.com/matteogiorgi/nn-option-pricing/blob/main/.github/workflows/docs.yml}{\path{.github/workflows/docs.yml}}.

At each push to the main branch, the workflow performs the following operations:

\begin{enumerate}
    \item checks out the repository;
    \item installs Python and the project dependencies;
    \item installs the Sphinx dependencies;
    \item builds the HTML site using the strict \texttt{-W} option;
    \item uploads the generated HTML site as a GitHub Pages artifact;
    \item deploys the artifact to GitHub Pages.
\end{enumerate}
This approach keeps the repository clean while still providing an online version of the API and project notes.
Once GitHub Pages is enabled for the repository, the site is available through the repository Pages URL.

In summary, the complete project can be inspected through two complementary online resources.
The repository contains the source code, report sources, scripts, configuration files, and experiment artifacts,
while the published Sphinx site provides a navigable view of the Python package, API reference, and project workflow:

\begin{itemize}
    \item \href{https://github.com/matteogiorgi/nn-option-pricing}{\texttt{GitHub repository}}.
    \item \href{https://matteogiorgi.github.io/nn-option-pricing}{\texttt{Online documentation}}.
\end{itemize}
